\documentclass[11pt,a4paper]{article}
\usepackage[margin=22mm]{geometry}
\usepackage[T1]{fontenc}
\usepackage{lmodern}
\usepackage{amsmath,amssymb}
\usepackage{microtype}
\usepackage{parskip}
\usepackage[hidelinks]{hyperref}
\hypersetup{pdftitle={Recovering Innovation Covariance}}

\begin{document}
\section*{Recovering Innovation Covariance}

Consider a stationary, causal VARMA model
$$
x_t=\sum_{i=1}^{p}A_i x_{t-i}
    +\sum_{j=0}^{q}B_j\varepsilon_{t-j},
\qquad B_0=I,
$$
where $x_t$ is $r$-dimensional and the innovations are temporally uncorrelated
with covariance $\Sigma$.

Suppose the coefficient matrices and the desired contemporaneous covariance
$$
V=\operatorname{Cov}(x_t,x_t)
$$
are given. We wish to determine $\Sigma$.

\subsection*{The Existing VYW System}

For $p\geq1$, we use the reduced vector Yule--Walker (VYW) system derived in
Appendix~B of J\'onasson and Ferrando (2006)~\cite{jf2006}. Their lagged covariance
matrices $S_k$ are denoted here by
$$
\Gamma_k=\operatorname{Cov}(x_t,x_{t-k}),
\qquad \Gamma_0=V,\qquad \Gamma_{-k}=\Gamma_k^T.
$$
The appendix eliminates $\Gamma_p$ and exploits the symmetry of $\Gamma_0$.
Let $\operatorname{vec}$ stack matrix columns and $\operatorname{vech}$ stack the
lower triangle column by column. Using the appendix's packing convention, define
$$
\begin{aligned}
v&=\operatorname{vech}(\widehat{\Gamma}_0),&
\sigma&=\operatorname{vech}(\Sigma),\\
h&=\operatorname{col}\bigl(
  \operatorname{vec}\Gamma_1,\ldots,\operatorname{vec}\Gamma_{p-1}\bigr).
\end{aligned}
$$
Here $\widehat{\Gamma}_0$ is lower triangular, with the same strictly lower
triangle as $\Gamma_0$ and half its diagonal, so that
$\Gamma_0=\widehat{\Gamma}_0+\widehat{\Gamma}_0^T$. The diagonal of $\Sigma$ is
\emph{not} halved in $\sigma$. The operator $\operatorname{col}$ stacks vectors;
$h$ is empty when $p=1$.

The reduced system can be written as
$$
F_0v+F_1h=g(\Sigma).
$$
The matrix $F=[\,F_0\;\;F_1\,]$ is the existing VYW coefficient matrix, constructed
as in Appendix~B, equations (B.3)--(B.5). It depends only on the AR coefficients.
Writing
$$
m=\frac{r(r+1)}{2},\qquad K=m+(p-1)r^2,
$$
$F_0$ has size $K\times m$ and $F_1$ has size $K\times(K-m)$.

\subsection*{Linearity of the Right-Hand Side}

The existing covariance recurrences are
$$
C_0=\Sigma,\qquad
C_j=B_j\Sigma+\sum_{i=1}^{\min(p,j)}A_iC_{j-i},
\quad j=1,\ldots,q,
$$
and
$$
G_k=\sum_{j=k}^{q}B_jC_{j-k}^T,\quad k=0,\ldots,q,
\qquad G_k=0\quad(k>q).
$$
Thus every $C_j$ and $G_k$ is linear in $\Sigma$. The reduced right-hand side
from Appendix~B is
$$
g(\Sigma)=\operatorname{col}\bigl(
  \operatorname{vech}(G_0^T+G_pA_p^T),
  \operatorname{vec}G_1,\ldots,\operatorname{vec}G_{p-1}\bigr)
=R\sigma,
$$
where $R$ is a $K\times m$ matrix determined by the AR and MA coefficients.
The $\operatorname{vech}$ operation in $g$ simply selects lower-triangle entries;
it does not halve diagonal entries.

\newpage
\subsection*{Forward and Inverse Computations}

In the usual forward computation, $\Sigma$ is known, and the VYW solve is
$$
\boxed{
\begin{bmatrix}F_0&F_1\end{bmatrix}
\begin{bmatrix}v\\h\end{bmatrix}=R\sigma.
}
$$
For the inverse computation, $\Gamma_0=V$ is known instead. Move its contribution
to the right-hand side and the contribution of $\Sigma$ to the left:
$$
\boxed{
\begin{bmatrix}R&-F_1\end{bmatrix}
\begin{bmatrix}\sigma\\h\end{bmatrix}=F_0v.
}
$$
The two systems have the same dimension $K$. Both $\Gamma_0$ and $\Sigma$ have
$m$ independent entries, so the inverse formulation exchanges which covariance
is prescribed while retaining the auxiliary lagged covariances as unknowns.
Only $\Sigma$ needs to be returned.

No new derivation of the Kronecker-product matrices is required: $F_0$ and $F_1$
come from the existing VYW construction. To assemble $R$, use the symmetric
basis matrices $E_a$ satisfying $\operatorname{vech}(E_a)=e_a$, where $e_a$ is
the $a$th coordinate vector in $\mathbb{R}^m$. Then
$$
R_{\cdot a}=g(E_a),\qquad a=1,\ldots,m.
$$
A diagonal basis matrix has one unit diagonal entry; an off-diagonal basis
matrix has a pair of symmetric unit entries. Each column requires only the
existing right-hand-side construction, not a full covariance solve.

This is an \emph{inverse VYW formulation}, not an adjoint or transpose problem.
The coefficient matrix changes, so a new LU factorization is needed.

\subsection*{Special Cases}

For a pure VAR(1) model, $\Gamma_1=A_1V$, giving the explicit formula
$$
\boxed{\Sigma=V-A_1VA_1^T.}
$$
For a pure VMA($q$) model, the equation is
$$
\boxed{V=\Sigma+\sum_{j=1}^{q}B_j\Sigma B_j^T.}
$$

\subsection*{Existence and Validity}

A positive-semidefinite target $V$ does not necessarily yield a positive-semidefinite
$\Sigma$. Thus not every target covariance is compatible with fixed VARMA
coefficients.

The inverse calculation must check for singularity or poor conditioning and verify
that the recovered $\Sigma$ is positive semidefinite and reproduces $V$ through
the forward covariance calculation.

\begin{thebibliography}{1}
\bibitem{jf2006}
J\'onasson, K. and Ferrando, S.~E. (2006).
\emph{Efficient likelihood evaluation for VARMA processes with missing values}.
Report VHI-01-2006, Faculty of Engineering, University of Iceland.
Appendix~B: \emph{Solution of the vector Yule--Walker equations}.
\end{thebibliography}

\end{document}
